\chapter{Vehicle Space Relative to the Railway State}
\label{chap:vehicle-space}

Vehicle Space does not introduce another railway state. It asks how
vehicle-related quantities are expressed relative to the qualified pose3
frame inherited from Part V and the traversal \(s(t)\) introduced here.

\begin{principle}[Track--Vehicle Distinction]
The qualified railway state may exist without a vehicle model. A vehicle
state requires that reference plus an identified vehicle model and operating
inputs.
\end{principle}

\begin{proposition}[State hierarchy]
\label{prop:state-hierarchy}
\[
\mathrm{vehicleState}_M(t)
=
\mathcal V_M\bigl(
\mathcal S_{\mathrm{rail}}(s(t)),Q_M
\bigr),
\]
but the railway state is not derived from
\(\mathrm{vehicleState}_M\). The model index and inputs remain part of the
result's meaning.
\end{proposition}

\section{Relative Quantities}

The moving track frame supplies longitudinal, lateral, and vertical reference
directions. A vehicle model may then express body or subsystem displacement,
yaw, roll, acceleration, wheel load, or envelope position relative to that
frame. These are examples of possible model outputs, not universally present
fields and not validated merely because they are expressed in Vehicle Space.

\[
F_{\mathrm{vehicle}}(t)\neq
F_{\mathrm{rail}}(s(t)).
\]
Their relative orientation is meaningful only after both frames, sign
conventions, time relation, and relevant vehicle degrees of freedom are
specified.

\section{Subsystem and Reference-Point Uses}

A wheelset, bogie, vehicle body, or selected centre-of-gravity point may
require a different state vector and model. In particular, a
centre-of-gravity trajectory is not the whole vehicle state. It may become an
explicit optimization output only if the vehicle model, objective, constraints
and applicability are declared.

\begin{summary}
Vehicle Space is a relative interpretation domain downstream of the qualified
railway state and traversal. It does not redefine pose3, the track frame, or
Alignment identity.
\end{summary}
